\documentclass[12pt]{article}
\usepackage{amsmath, amsfonts, amssymb, amsthm, fullpage, color, bbm, enumerate}
\usepackage[T1]{fontenc}
\usepackage{lmodern}

\usepackage[onehalfspacing]{setspace}

\usepackage[hidelinks]{hyperref}
%\usepackage[capitalize,nameinlink,noabbrev]{cleveref}

\theoremstyle{plain}
\newtheorem{asn}{Assumption}
%\crefname{asn}{Assumption}{Assumptions}
\newtheorem{lem}{Lemma}
%\crefname{lem}{Lemma}{Lemmas}
\newtheorem{thm}{Theorem}
\newtheorem{prop}{Proposition}
\newtheorem{cor}{Corollary}
\newtheorem{claim}{Claim}

\DeclareMathOperator*{\argmax}{arg\,max}
\DeclareMathOperator*{\argmin}{arg\,min}
\DeclareMathOperator*{\tr}{tr}
\DeclareMathOperator*{\rk}{rk}
\DeclareMathOperator*{\spa}{span}
\DeclareMathOperator*{\ve}{vec}
\DeclareMathOperator*{\diag}{diag}
\DeclareMathOperator*{\Rea}{Re}
\DeclareMathOperator*{\var}{Var}
\DeclareMathOperator*{\cov}{Cov}
\DeclareMathOperator*{\corr}{Corr}

\setlength{\skip\footins}{\baselineskip}
\setlength{\footnotesep}{0.3\baselineskip}

\allowdisplaybreaks

\begin{document}

{\noindent \bf University of Chicago \hfill
Spring 2026 \\
Department of Economics \hfill ECMA 31150}

\begin{center}
{\Large Problem Set 2} \\
This version: \today \\[1ex]
Mikkel Plagborg-Moller \\
\vspace{0.5\baselineskip}
Due Friday, April 17
\end{center}
\vspace{0.2\baselineskip}

\noindent \emph{Please upload your answers on Canvas, with code attached. You are encouraged to collaborate, but please write up your answers independently and do not share code chunks. Use any software you like, although you learn the most by limiting the use of generative AI. If you find any errors or require clarification, please let me know right away.}

\subsection*{Question 1}
Consider a bivariate VAR(2) model for $(y_t,x_t)$:
\begin{align*}
y_t &= 1 + 0.7 y_{t-1} + 0.1 y_{t-2} + 0.2 x_{t-1} - 0.2x_{t-2} + u_{1t}, \\
x_t &= 0.6 x_{t-1} + 0.3 x_{t-2} + 0.2 y_{t-2} + u_{2t},
\end{align*}
where $u_{1t}$ and $u_{2t}$ are i.i.d.\ bivariate normally distributed, each with variance 1, and $\corr(u_{1t},u_{2t})=0.4$.
\begin{enumerate}[i)]
\item Is the model covariance stationary?
\item Does $\lbrace x_t \rbrace$ Granger-cause $\lbrace y_t \rbrace$? Does $\lbrace y_t \rbrace$ Granger-cause $\lbrace x_t \rbrace$?
\item What is the MSFE of the optimal one-step-ahead forecast of $y$, i.e., $E[y_{t+1} \mid \lbrace y_s,x_s \rbrace_{s \leq t}]$? What is the MSFE of the optimal two-step-ahead forecast of $y$, i.e., $E[y_{t+2} \mid \lbrace y_s,x_s \rbrace_{s \leq t}]$?
\item What is $E(y_t)$, $E(x_t)$, $\var(y_t)$, $\var(x_t)$, and $\corr(y_t,x_t)$?
\item Simulate a very large data set with $T=10000$ from this VAR model. Use the simulation output to (approximately) verify your calculations in part (iv).
\end{enumerate}
[Hint: some of these questions will be hard to answer with ``pen and paper'', so feel free to solve them numerically on the computer.]

\subsection*{Question 2}
Consider a univariate AR($p$) model
\[y_t = c + a_1y_{t-1} + a_2 y_{t-2} + \dots + a_p y_{t-p} + u_t,\quad u_t \sim \mathit{WN}(\sigma_u^2).\]
Let $z_1,\dots,z_p$ denote the $p$ (possibly complex) roots of the characteristic equation
\[1 = a_1 z + a_2 z^2 + \dots + a_p z^p,\]
and assume these are all distinct. Finally, let $\tilde{\boldsymbol{A}}$ be the $p \times p$ companion matrix that re-expresses the univariate AR($p$) model as a multivariate VAR(1) model.
\begin{enumerate}[i)]
\item Show that the eigenvalues and corresponding eigenvectors of $\tilde{\boldsymbol{A}}$ are given by
\[z_i^{-1}\quad \text{and} \quad (1,z_i,z_i^2,\dots,z_i^{p-1})',\]
respectively, for $i=1,\dots,p$.
\item How does the result in part (i) relate to the methods for checking covariance stationarity that we have discussed in lecture? Be brief.
\end{enumerate}

\subsection*{Question 3}
Go to the St.\ Louis Fed's FRED website and download data on real GDP (FRED code: GDPC1), the consumer price index (code: CPIAUCSL), and consumer sentiment (code: UMCSENT). Collapse the last two series to quarterly averages.\footnote{Note that UMCSENT is recorded at a quarterly frequency until 1978.} Then transform real GDP and the consumer price index to quarterly log growth rates. Use only data from the beginning of 1960 until (and including) the end of 2019. Consider a VAR model in the three variables (real GDP growth, CPI inflation, consumer sentiment).

\begin{enumerate}[i)]
\item What lag length is selected by the AIC for the VAR? Use a maximum potential lag length of 12 quarters.
\end{enumerate}
Regardless of the result from part (i), we henceforth assume a VAR(1) model.
\begin{enumerate}[i)]
\setcounter{enumi}{1}
\item Estimate the model. Is the estimated model covariance stationary?
\item Assume (for this part only) that the VAR innovations are i.i.d.\ Gaussian. Compute an 80\% forecast interval for inflation in the first quarter of 2020. [Hint: we did not explicitly discuss how to construct VAR forecast intervals in lecture, but you should be able to adapt the arguments from the slides on AR models.]
\item Compute 90\% bootstrap confidence intervals for each of the coefficients in the equation for inflation. Report these confidence intervals alongside the point estimates in a nicely formatted table. Use at least 500 bootstrap iterations.
\end{enumerate}


\end{document}
